\documentclass[12pt,a4paper]{article}
\usepackage[utf8]{inputenc}
\usepackage[brazil]{babel}
\usepackage{graphicx}
\usepackage{amssymb, amsfonts, amsmath}
\usepackage{float}
\usepackage{enumerate}
\usepackage[top=2.5cm, bottom=2.5cm, left=1.25cm, right=1.25cm]{geometry}

\begin{document}
\pagestyle{empty}

\begin{center}
  \begin{tabular}{ccc}
    \begin{tabular}{c}
      \includegraphics[scale=0.25]{../../biblioteca/imagem/brasao-de-armas-brasil} \\
    \end{tabular} & 
    \begin{tabular}{c}
      Ministério da Educação \\
      Universidade Federal dos Vales do Jequitinhonha e Mucuri \\
      Faculdade de Ciências Sociais, Aplicadas e Exatas - FACSAE \\
      Departamento de Ciências Exatas - DCEX \\
      Disciplina: Matemática I \quad Semestre: 2020/1\\
      Prof. Me. Luiz C. M. de Aquino\\
    \end{tabular} &
    \begin{tabular}{c}
      \includegraphics[scale=0.25]{../../biblioteca/imagem/logo-ufvjm} \\
    \end{tabular}
  \end{tabular}
\end{center}

\begin{center}
  \textbf{Lista I}
\end{center}

\begin{enumerate}
  \item Considere os conjuntos $A = \left\{-5,\,-4,\,8,\,10\right\}$,
    $B = \left\{x\in\mathbb{Z} \,|\, x < 8\right\}$ e
    $C = \left\{x\in\mathbb{N} \,|\, x > 6\right\}$.

    Determine:
    \begin{enumerate}
      \item $A\cup B$
      \item $B\cap C$
      \item $B\setminus A$
      \item $(B\cup C)\cap A$
    \end{enumerate}

  \item Suponha que $A = \{x\in\mathbb{N} \,|\, x^2 - x - 2 = 0\}$ e 
    $B = \{x\in\mathbb{Z} \,|\, -5 \leq x < 4 \}$. Classifique em Verdadeiro ou
    Falso:

    \begin{enumerate}[(\ \ )]
      \item $A\cap B = \{-1,\,2\}$
      \item $n(B) = 10$
      \item $A\cup B = A$
      \item $A=\varnothing$
      \item $B - A=\varnothing$
    \end{enumerate}

  \item Classifique cada afirmação abaixo como Verdadeiro ou Falso.

    \begin{enumerate}[(\ \ )]
      \item $-3 \in \{x\in \mathbb{N} \,|\, x < 4\}$
      \item Se $a$, $b\in\mathbb{N}$, então $a - b \in\mathbb{N}$.
      \item Se $a$, $b\in\mathbb{R} - \mathbb{Q}$, então 
        $a + b \in\mathbb{R} - \mathbb{Q}$.
      \item Se $a\notin \mathbb{Q}$, então $a\neq 0$.
      \item $\{1\} \subset \{\{1\},\,2,\,3\}$
      \item Se $A = \{1,\,2\}$ e $B = \{1,\,2,\,3\}$, então $A\in B$.
      \item Se $A - B = \varnothing$ e $B - A = \varnothing$, então $A = B$.
      \item Se $ab \in \mathbb{N}$, então $a \in \mathbb{N}$ e $b \in \mathbb{N}$.
      \item Se $A\cap B = \varnothing$, então $n(A\cup B) = n(A) + n(B)$.
      \item Se $n\in A$ e $n\in B$, então $n\in A\cap B$.
    \end{enumerate}

  \item Considerando o Diagrama de Venn padrão para três conjuntos $A$, $B$ e
  $C$, represente cada operação abaixo com seu respectivo diagrama.

    \begin{enumerate}[(a)]
      \item $A\cap(B\cup C)$
      \item $(A\cup B)\setminus C$
      \item $(A\setminus B)\cup(C\setminus B)$
      \item $(A\cup B)\cap(A\cup C)$
    \end{enumerate}

\end{enumerate}

\end{document}